\documentclass[11pt,letterpaper]{article}

\usepackage[margin=1in]{geometry}
\usepackage{times}
\usepackage{amsmath,amssymb}
\usepackage{setspace}
\usepackage[hidelinks]{hyperref}
\usepackage{microtype}

\pagestyle{empty}
\setstretch{1.08}
\setlength{\parskip}{0.55em}
\setlength{\parindent}{0pt}

\newcommand{\abssection}[1]{\textbf{#1.}~}

\begin{document}

\begin{center}
{\LARGE\bfseries Suppressed but Unmoved:}\\[0.35em]
{\large\bfseries Defensive Suppression and Failed Offensive Adaptation in the NBA}\\[1em]
{\large Rami Zheman}\\[0.25em]
{\normalsize\textit{Independent Sports Analytics Researcher}\\
Alexandria, VA}
\end{center}

\vspace{0.6em}

\abssection{Introduction}
A coach watches a play that was working suddenly stop working.
Execution? Random variation? Or has the defense taken it away?
Coaches talk constantly about adjustments, yet NBA offenses rarely
make them. More than nine times in ten, teams keep calling the same
play at nearly the same rate after it breaks down. This paper argues
that in many cases the defense has already won, and continuing the
play is a costly mistake. This matters for the industry because if
such collapses are real and go unaddressed, teams are forfeiting
points in real time using information that is already sitting in
front of them mid-game.

\abssection{Methods}
The analysis draws on 863{,}752 classified NBA possessions from the
2022--23 through 2025--26 seasons (5{,}234 games), built from public
play-by-play through the NBA Stats API, with a calibrated play-type
classifier covering pick-and-roll, spot-up, drive, post-up, cut,
pull-up, transition, and putback. A play type is flagged as a
suppression trigger when a team's first-half points per possession
exceeds its second-half figure by $0.25$ or more, with at least five
possessions in each half. Because triggers are selected on an elevated
first half, a raw before-and-after comparison conflates suppression
with ordinary mean reversion. This is corrected with a same-game
matched-pair design: each triggered play type is paired with the
opposing team running the same play type in the same game, matched on
first-half efficiency, with cluster-robust and permutation inference,
plus caliper, blowout, and overtime-boundary robustness checks.

\abssection{Results}
Defensive suppression holds up net of regression to the mean, with
matched pairs showing an excess decline of $0.408$ points per
possession (95\% CI $0.378$ to $0.438$, $N = 1{,}136$ pairs across
848 games, $p < 0.0002$), stable across every robustness check.
Offenses largely fail to adapt: teams continue running the suppressed
play in 93\% of events, at 97 to 98\% of prior volume, at a cost of
roughly $5.4$ to $5.5$ points per trigger. Substituting the primary
ball-handler does not recover efficiency (a negligible $0.039$ point
difference). Even when a higher-scoring alternative already existed in
the same game, true of 95.3\% of triggers, no team systematically
reallocated toward it. Matchup-ranked counters, validated out of
sample, outperform continuation by $0.445$ points per possession in
the playoffs and $0.584$ in the regular season, while team-level
suppression severity and unrealized counter value show only a weak
relationship with win rate ($r = +0.25$).

\abssection{Conclusion}
Defensive suppression is a genuine, sizable, and largely unaddressed
source of offensive inefficiency. Teams continue running actions the
defense has already solved, even when better options are available in
the same game, and a simple first-half signal can identify both the
collapse and its remedy. This suggests a practical application for
the industry: real-time tools that flag suppression as it happens and
rank live counters. Whether teams acting on the signal can convert it
into realized efficiency remains an open question for future work.

\vspace{0.8em}
\abssection{AI Disclosure}
Large language models (Claude, Anthropic; Grok, xAI) were used to
help write the analysis code and draft this paper's text. The author
directed the research question, supplied and validated the data, and
reviewed and edited all results and text.

\end{document}
